
Our work sits at the intersection of three research streams: the theory of
asymmetric loss and quantile regression, the literature on probabilistic
electricity forecasting, and the emerging body of work on cost-oriented learning
for energy applications.

\subsection{Asymmetric Loss and Quantile Regression}

The theoretical foundation for prediction under asymmetric costs was laid by
\textcite{koenker1978regression}, who introduced
\emph{regression quantiles} as a natural generalization of sample quantiles to
the linear model. They showed that minimizing the asymmetric pinball loss
(equivalently, the ``check function'') yields an estimator of the conditional
$\tau$-quantile of the response, and established its asymptotic properties.
This result directly underpins our use of a cost-sensitive asymmetric loss
function --- instantiated as a non-linear variant of the pinball loss --- as
the training objective for day-ahead demand forecasting.

\textcite{christoffersen1997optimal} provided the
corresponding decision-theoretic foundation. They showed that under a general
asymmetric loss function, the optimal point forecast is \emph{biased}, and that
the optimal amount of bias is time-varying and depends on higher-order
conditional moments such as conditional variance. In particular, their analysis
implies that any forecast trained with a symmetric objective (e.g., MSE) will be
systematically suboptimal when the true cost structure is asymmetric, even if
the conditional mean is estimated perfectly.

\subsection{Electricity Forecasting: Overview and Probabilistic Methods}

\textcite{weron2014electricity} provides a comprehensive
review of the state of the art in electricity price forecasting, surveying
statistical, machine-learning, and hybrid approaches. While their focus is on
price rather than demand, the review highlights the central role of model
selection, loss function choice, and benchmark comparisons that also apply in
the load forecasting context.

\textcite{hong2016probabilistic} offer a tutorial
review of probabilistic electric load forecasting, covering methods ranging from
regression-based quantile estimation to ensemble and Bayesian approaches.
Crucially, they argue that point forecasts are insufficient for operational
decision-making and advocate for full predictive distributions or at least
prediction intervals — a view consistent with the quantile-forecasting objective
we adopt.

The practical importance of probabilistic evaluation metrics, and of the pinball
loss in particular, was cemented by the Global Energy Forecasting Competition
2014~\cite{hong2016gefcom2014}. GEFCom2014 was the first major forecasting
competition to use pinball loss as its primary scoring rule across electricity
load, price, wind, and solar tracks. This competition established pinball loss
as a community standard and demonstrated that models explicitly optimized for
quantile accuracy substantially outperform those evaluated post-hoc by quantile
metrics.

\subsection{Cost-Oriented and Asymmetric Learning for Load Forecasting}

The most directly related line of work addresses the mismatch between symmetric
training objectives and asymmetric operational costs in electricity markets.

\textcite{wang2019pinball} train a Long Short-Term Memory
(LSTM) network directly with the pinball loss to produce probabilistic individual
load forecasts. By guiding the network to minimize $\mathbb{E}[\mathcal{L}_\tau]$
rather than MSE, their approach yields calibrated quantile predictions without
requiring post-hoc quantile regression on the residuals. Their results show
improved reliability and sharpness compared to symmetric-loss baselines, and the
architecture serves as a direct precursor to our approach.

\textcite{lin2020asymmetric} address the related problem of
contract capacity optimization for high-voltage consumers, where the electricity
pricing structure is inherently asymmetric: underestimating demand triggers
penalty charges at a higher per-unit rate than overestimating. They propose
several asymmetric loss functions derived directly from the Taiwan Power
Company's pricing formula, and demonstrate that LSTM models trained with these
loss functions reduce electricity costs by 0.88\%--2.42\% over symmetric MSE
baselines. Their work shows that embedding domain-specific cost asymmetry into
the learning objective is preferable to symmetrically minimizing forecast error.

\textcite{wu2021svr} extend this idea to support vector regression
(SVR), developing a cost-oriented asymmetric SVR (AsySVR) framework based on a
linear-linear loss function with an insensitivity parameter to improve
generalization. Evaluated on New South Wales load data, their asymmetric
framework achieves daily economic cost reductions of 42\%--57\% relative to
standard SVR, underscoring that the choice of loss function can have a large
practical impact on operational outcomes.

\textcite{zhang2022costoriented} further formalize the
decision-oriented perspective on load forecasting, framing the problem explicitly
in terms of the downstream cost incurred by the system operator. Their
cost-oriented framework bridges the gap between forecast accuracy metrics and
operational objectives, and provides theoretical justification for replacing
standard point-forecast losses with cost-aware alternatives.

Taken together, this body of work establishes that asymmetric, cost-sensitive
training objectives are both theoretically motivated and empirically superior to
symmetric alternatives in electricity applications. Our work builds on these
foundations by applying a cost-sensitive asymmetric loss function --- instantiated
as a non-linear variant of the pinball loss --- within the newsvendor framework,
directly targeting the cost-minimising quantile identified in
Section~\ref{sec:modeling}.
